\section{Linguistic Note: “Dream” as Vision and Aspiration}\label{sec:linguistic_note}

In many languages, the same word denotes both the visions of sleep and the aspirations of waking life—Portuguese \textit{sonho}, Spanish \textit{sueño}, Japanese \textit{yume}, Chinese \textit{mèng}, and others.  

This semantic overlap helps explain why dreams have so often been treated as sources of direction or vocation rather than mere nocturnal noise.  
Throughout this work, when “dream” refers to an \emph{aspiration}, the context will make this explicit to avoid ambiguity.

\paragraph{To be developed in~\ref{part:science_of_sleep}:}
Later sections will connect these cultural practices to specific neurophysiological sleep states (N1/N2/REM) and examine how each ancient intuition corresponds (or fails to correspond to) contemporary data.